% Trinity S³AI — Boundary-Mapping Research Program in Geometric Unification
% Wave 17.3 — arXiv revision (Waves 16–17 results incorporated)
\documentclass[11pt,a4paper]{article}
\usepackage[utf8]{inputenc}
\usepackage[T1]{fontenc}
\usepackage{amsmath,amssymb,amsthm,mathtools}
\usepackage{hyperref}
\usepackage{booktabs}
\usepackage{slashed}
\usepackage{graphicx}
\usepackage{tikz}
\usepackage[numbers,sort&compress]{natbib}

% Theorem environments
\theoremstyle{definition}
\newtheorem{theorem}{Theorem}[section]
\newtheorem{lemma}[theorem]{Lemma}
\newtheorem{definition}[theorem]{Definition}
\newtheorem{corollary}[theorem]{Corollary}
\newtheorem{conjecture}[theorem]{Conjecture}
\newtheorem{boundary}{Boundary Theorem}

% Macros
\newcommand{\ph}{\varphi}
\newcommand{\Dirac}{\slashed{D}}
\newcommand{\tr}{\operatorname{tr}}
\newcommand{\spec}{\operatorname{Spec}}
\newcommand{\Ad}{\operatorname{Ad}}
\newcommand{\Aut}{\operatorname{Aut}}
\newcommand{\sgn}{\operatorname{sgn}}
\newcommand{\R}{\mathbb{R}}
\newcommand{\C}{\mathbb{C}}
\newcommand{\Z}{\mathbb{Z}}
\newcommand{\N}{\mathbb{N}}

\title{Trinity S$^3$AI: An Active Boundary-Mapping Research Program in Geometric Unification}
\author{[Author List TBD]\\
  \texttt{trinity-s3ai@proton.me}}
\date{May 2026}

\begin{document}
\maketitle

% ============================================================
\begin{abstract}
We present a systematic test of the hypothesis that the parameters of the
Standard Model (SM) derive from the geometry of the H$_4$ Coxeter group and
the 600-cell. Using a combination of formal proofs (Coq/Rocq), numerical
validation (Python), and discrete spectral geometry, we evaluate every
verifiable claim of the ``Trinity'' program.

Our \textbf{positive results} include: (i) the eta-invariant
$\eta(S^3/2I) = -2$ via the Atiyah--Patodi--Singer theorem applied to the
E$_8$ plumbing; (ii) KO-dimension $6 \pmod 8$ for the H$_4$/600-cell Dirac
operator, matching the Connes--Chamseddine SM requirement; (iii) numerical
verification of 25 SM parameter formulas against PDG 2024 with mean relative
error $0.1\%$; (iv) JUNO 2025 and DESI DR2 data consistent with Trinity
oscillation and cosmology predictions.

Our \textbf{boundary theorems} include seven Boundary theorems:
$\sigma$-field unavailable in H$_4$ (BT-1); no 3-generation structure from
H$_4$ (BT-2); D$_4$ KO-dimension $5 \pmod 8 \neq 6$ (BT-3); spectral
action Higgs mass $132.88\,\text{GeV} \neq 125.10\,\text{GeV}$ (BT-4);
E$_6$/E$_7$ plumbing discretization mismatch (BT-5); F$_4$ Yukawa scan
fails to produce 3-generation hierarchy without external parameters (BT-7).
The formal proof base now stands at \textbf{$1375$~Qed}, \textbf{$0$~Admitted},
across \textbf{$79$~Coq files} and \textbf{$77$~Axioms}, with an active
Lean~4 port.

We discuss the surviving open directions and conclude that naive geometric
unification via discrete Coxeter groups is \textbf{refuted} as a complete SM
derivation, while retaining partial structural insights.
\end{abstract}

% ============================================================
\section{Introduction}\label{sec:intro}

The idea that the Standard Model (SM) Lagrangian encodes hidden geometric
structure dates back to the Kaluza--Klein program and finds its modern
expression in Connes' noncommutative geometry (NCG)
\cite{Connes:1996,ConnesMarcolli:2008}. In NCG, the SM gauge group, fermion
representations, and Higgs mechanism emerge from the spectral triple
$(\mathcal{A}, \mathcal{H}, \Dirac)$ of a suitable ``internal'' space.

The ``Trinity S$^3$AI'' program (Waves 1--17) tests a specific incarnation
of this idea: the internal space is the quotient $S^3/2I$ of the 3-sphere by
the binary icosahedral group $2I$, discretized via the 600-cell (the
H$_4$ root polytope). The conjecture is that the 120 vertices of the
600-cell, equipped with a discrete Dirac operator and spectral action,
reproduce the SM parameters without free fitting.

\textbf{Scope and honesty.} This paper reports a \emph{boundary-mapping research program}. We do \emph{not} claim to disprove NCG in general; we refute a
\emph{specific discrete model}. Every theorem is either formally proved in
Coq (\S\ref{sec:coqstats}), numerically verified in Python, or honestly
labelled as open.

\textbf{Roadmap.} \S\ref{sec:framework} reviews the discrete spectral
framework. \S\ref{sec:positive} states the surviving positive results.
\S\ref{sec:boundary} presents the five original Boundary theorems.
\S\ref{sec:extended} reports new boundary findings from Waves 16--17
(F$_4$ Yukawa scan and E$_8$ plumbing convergence).
\S\ref{sec:discussion} discusses implications, and
\S\ref{sec:conclusion} summarises open questions.

% ============================================================
\section{Mathematical Framework}\label{sec:framework}

\subsection{Discrete Dirac operators on $S^3/\Gamma$}

Let $\Gamma \subset SU(2)$ be a finite subgroup (binary polyhedral). The
quotient $M = S^3/\Gamma$ is a spherical space form. A discrete Dirac
operator is defined on the vertex set $V$ of a $\Gamma$-invariant polytope:
\begin{equation}\label{eq:dirac}
  (D\psi)_i = \sum_{j \sim i} \langle i j \rangle \cdot \gamma(\overrightarrow{ij}) \, \psi_j,
\end{equation}
where $j \sim i$ denotes nearest neighbours, $\gamma$ is the Clifford
representation, and $\langle i j \rangle$ are coupling constants.

\subsection{KO-dimension}

The real structure $J$ (charge conjugation) must satisfy Connes' axioms
\cite{Connes:1996}:
\begin{equation}
  J^2 = \varepsilon, \quad JD = \varepsilon' DJ, \quad J\Gamma = \varepsilon'' \Gamma J.
\end{equation}
For the SM one needs KO-dimension $6 \pmod 8$, i.e.
$(\varepsilon, \varepsilon', \varepsilon'') = (+, +, -)$.

\subsection{Spectral action}

The bosonic action is
\begin{equation}\label{eq:spectral-action}
  S_\Lambda(D) = \tr\bigl(f(D/\Lambda)\bigr),
\end{equation}
with $f$ a smooth cutoff. The heat-kernel expansion yields effective
couplings. In the Connes--Chamseddine model the tree-level Higgs mass is
$m_H^2 \propto \Lambda^2$, brought down to the observed $125\,\text{GeV}$
by a $\sigma$-field correction \cite{ConnesMarcolli:2008}.

\subsection{Formal proof infrastructure}\label{sec:coqstats}

All Coq proofs use the Rocq/Coq~8.20.1 proof assistant
with the \texttt{coq-interval} library for automated real bounds.
As of Wave~17:
\begin{itemize}
  \item $1375$ theorems with \texttt{Qed}
  \item $0$ theorems remain \texttt{Admitted}
  \item $77$ Axioms (all tagged with physical or mathematical justification)
  \item $79$ \texttt{.v} files in the canonical proof tree
  \item Lean 4 port: $9$ modules, $1$ honest \texttt{sorry} (needs Mathlib)
\end{itemize}

% ============================================================
\section{Positive Results}\label{sec:positive}

\begin{theorem}[Eta-invariant, Wave 8.3]\label{thm:eta}
  The eta-invariant of the Dirac operator on $S^3/2I$ is
  \begin{equation}
    \eta(D_{S^3/2I}) = -2.
  \end{equation}
\end{theorem}
\begin{proof}[Sketch]
  $S^3/2I = \Sigma(2,3,5)$ bounds the E$_8$ plumbing $P_{E_8}$.
  By APS:
  $0 = \int_{P_{E_8}} \hat{A} - (\eta + h)/2$.
  With $\hat{A}(P_{E_8}) = -1$ and $h = 0$ one obtains $\eta = -2$.
  Full proof in \texttt{EtaDFBridge.v}.
\end{proof}

\begin{theorem}[KO-dimension, Wave 5.1]\label{thm:ko}
  The discrete real structure $J$ on the H$_4$/600-cell spinor space
  satisfies
  \begin{equation}
    (\varepsilon, \varepsilon', \varepsilon'') = (+, +, -),
  \end{equation}
  i.e.\ KO-dimension $6 \pmod 8$.
\end{theorem}

\begin{theorem}[Formula verification, Wave 4]\label{thm:formulas}
  $25$ SM parameter formulas (mass ratios, mixing angles, gauge couplings)
  are verified against PDG 2024 with mean relative error $0.10\%$.
\end{theorem}

% ============================================================
\section{Boundary Theorems}\label{sec:boundary}

\begin{boundary}[Sigma-field, Wave 5.3]\label{boundary:sigma}
  The Hodge $*$ operator required for the Connes--Marcolli $\sigma$-field
correction does not exist on the 600-cell triangulation with the correct
sign $*^2 = -1$ on 2-forms.
\end{boundary}
\textit{Consequence:} Without $\sigma$, the spectral action predicts
$m_H \approx 132.88\,\text{GeV}$ (Wave 11.5), incompatible with the
observed $125.10 \pm 0.14\,\text{GeV}$ at $55.6\sigma$.

\begin{boundary}[Three generations, Wave 9]\label{boundary:3gen}
  The H$_4$ outer automorphism group is trivial ($|\operatorname{Out}(H_4)| = 1$).
  Therefore the 600-cell Dirac spectrum does not decompose into three
  isomorphic fermion families.
\end{boundary}

\begin{boundary}[D$_4$ KO-dimension, Wave 11.2]\label{boundary:d4}
  The D$_4$/24-cell discrete real structure has KO-dimension
  $5 \pmod 8$ ($J^2 = -1$), not $6 \pmod 8$.
\end{boundary}

\begin{boundary}[Higgs mass, Wave 11.5]\label{boundary:higgs}
  The Chamseddine--Connes spectral action on the 600-cell predicts a
  tree-level Higgs mass of $132.88\,\text{GeV}$, deviating from the PDG
  value $125.10\,\text{GeV}$ by $7.78\,\text{GeV}$ ($55.6\sigma$).
  One-loop corrections increase the discrepancy to $\sim 175\,\text{GeV}$
  (Wave 12.4).
\end{boundary}

\begin{boundary}[E$_6$/E$_7$ plumbing, Wave 12.5]\label{boundary:e6e7}
  Explicit discrete Dirac operators on the E$_6$ and E$_7$ plumbings yield
  boundary eta-invariants $-4$ and $-2$ respectively, which differ from the
  continuum targets $-3/2$ and $-7/4$ due to finite-size discretization.
\end{boundary}

% ============================================================
\section{Extended Boundary Findings (Waves 16--17)}\label{sec:extended}

\subsection{F$_4$ Yukawa Test}\label{sec:f4yukawa}

The F$_4$ root system (48 roots, binary octahedral group $2O$) was proposed
as an alternative geometric source of fermion masses because F$_4$ contains
D$_4$ as its long-root subalgebra and D$_4$ possesses triality
($\operatorname{Out}(D_4) \cong S_3$), a structure historically invoked for
three generations \cite{Ramond:2001,Gursey:1975}. We performed a systematic
numerical scan of Yukawa matrices built from F$_4$ root inner products.

\textbf{Scan design.} Five models were tested per fermion type (up, down,
lepton): (i) \textit{symmetric sector} --- equal averaging over the three
D$_4$ triality sectors; (ii) \textit{weighted sector} --- sector-dependent
weights; (iii) \textit{full 48} --- all F$_4$ roots with one coupling; (iv)
\textit{full 24 short} --- short roots only; (v) \textit{block 48} ---
long--short block matrix with four free parameters.

\textbf{Structural obstruction.} The F$_4$ short roots decompose as
$8_v \oplus 8_s \oplus 8_c$ under D$_4$ triality. Any symmetric averaging
over these three sectors yields a matrix of the form $aI + b(J-I)$, whose
eigenvalues are $\{a+2b, a-b, a-b\}$. Consequently, \emph{without explicit
symmetry breaking}, at most two distinct eigenvalues are possible. The
maximum hierarchy achievable by symmetric models is set by the geometry of
the 24-cell and is approximately $20{:}1$ --- five orders of magnitude short
of the observed up-quark hierarchy ($\sim 10^5$).

\textbf{Numerical results.} The symmetric-sector chi-squared values exceed
$21$ for all fermion types (Table~\ref{tab:f4}). The weighted and block
models can \emph{fit} the data, but require $3$--$4$ free parameters per
fermion type with \emph{no geometric origin} in F$_4$. They therefore play
the same role as the SM's own Yukawa couplings, defeating the purpose of a
derivation.

\begin{table}[htb]
\centering
\caption{F$_4$ Yukawa scan results (Wave 16.1).}
\label{tab:f4}
\begin{tabular}{lccc}
\toprule
Model & Up-type $\chi^2$ & Down-type $\chi^2$ & Lepton $\chi^2$ \\
\midrule
Symmetric sector & $32.5$ & $9.0$ & $28.1$ \\
Weighted sector & $4.7 \times 10^{-18}$ & $8.9 \times 10^{-19}$ & $2.4 \times 10^{-18}$ \\
Full 48 & $33.4$ & $9.0$ & $28.1$ \\
Full 24 short & $32.8$ & $9.0$ & $28.1$ \\
Block 48 ($4$ params) & $6.6 \times 10^{-18}$ & $8.3 \times 10^{-18}$ & $2.2 \times 10^{-20}$ \\
\bottomrule
\end{tabular}
\\[0.5em]
\small\textit{Note:} Low $\chi^2$ in weighted/block models comes from
fitting $3$--$4$ free parameters, not from geometric prediction.
\end{table}

\begin{boundary}[F$_4$ three-generation hierarchy, Wave 16.1]\label{boundary:f4}
  F$_4$ root geometry alone cannot produce the observed three-generation
  fermion mass hierarchy. The D$_4$ triality symmetry inherent in the
  short-root sector forces eigenvalue degeneracies that cap symmetric-model
  mass ratios at $\sim 20{:}1$. Models with enough free parameters can fit
  the data, but those parameters lack a principled geometric derivation.
\end{boundary}

\subsection{E$_8$ Plumbing Convergence}\label{sec:e8conv}

Theorem~\ref{thm:eta} states $\eta = -2$ via the APS theorem on the
E$_8$ plumbing. A natural question is whether the \emph{discrete} Dirac
operator on the plumbing converges to the same value as the discretization
is refined.

We constructed a $512 \times 512$ discrete Dirac operator $D_P$ for the full
E$_8$ plumbing (8 interior vertices + 120 boundary vertices = 600-cell) using
a geometrically motivated heuristic $B$ matrix that couples interior plumbing
nodes to icosidodecahedral equatorial slices of the boundary
\cite{ConwaySmith:2003}. The operator is Hermitian, $75\%$ sparse, and its
boundary block anticommutes with $\Gamma^5$.

\textbf{Eta results.} Sign-counting on the boundary block yields
$\eta_{\text{sign}} = -64$, far from the APS target $-2$. Soft-cutoff
regularisation
\begin{equation}
  \eta_{\text{reg}}(\varepsilon) = \sum_{\lambda \neq 0}
  \sgn(\lambda)\, e^{-\varepsilon|\lambda|}
\end{equation}
converges to $\eta_{\text{reg}} \approx -65.93$ as $\varepsilon \to 0$
(Table~\ref{tab:e8eta}). The reduced model (20 vertices) gave
$\eta = 0$; the full 128-vertex model gives $\eta \approx -66$. Neither
reproduces the continuum target.

\begin{table}[htb]
\centering
\caption{Discrete eta-invariants for the E$_8$ plumbing (Wave 16.2).}
\label{tab:e8eta}
\begin{tabular}{lcc}
\toprule
Model & $\eta_{\text{discrete}}$ & APS target \\
\midrule
Reduced ($20$ vertices, sign-count) & $0$ & $-2$ \\
Full ($128$ vertices, sign-count) & $-64$ & $-2$ \\
Full ($128$ vertices, $\varepsilon = 0.20$) & $-51.98$ & $-2$ \\
Full ($128$ vertices, $\varepsilon = 0.10$) & $-61.13$ & $-2$ \\
Full ($128$ vertices, $\varepsilon = 0.05$) & $-64.55$ & $-2$ \\
Full ($128$ vertices, $\varepsilon = 0.01$) & $-65.93$ & $-2$ \\
\bottomrule
\end{tabular}
\\[0.5em]
\small\textit{Note:} The B-matrix is a heuristic, not rigorously derived.
Agreement within $O(1)$ is expected for a finite discretisation.
\end{table}

\textbf{Honest assessment.} The discrete eta-invariant does \emph{not}
converge to $-2$ with the current heuristic $B$ matrix. Whether a
rigorously derived $B$ (from the plumbing geometry) would improve the
agreement remains an \textbf{open problem}. The sign-counting discrepancy
is understood: sign-counting is not zeta-regularisation, and the 600-cell
graph has abundant near-zero modes whose regularisation determines the
answer.

% ============================================================
\section{Discussion}\label{sec:discussion}

\subsection{Statistical assessment}\label{sec:stats}

We subjected the 25 core formulas to a pre\-registered Monte\-Carlo test
against the null hypothesis of random coincidence.  The search space
was fixed \emph{a priori}: $C\,\phi^a\pi^b e^c$ and four related templates,
with $C$ ranging over reduced rationals $p/q\le 1000$ and
$a,b,c\in\{-6,\dots,6\}$ (approx. $6.7\times10^8$ distinct formulas).
For each of $5\times10^4$ trials we drew $2\times10^3$ random formulas
and recorded the best relative error per observable.

\begin{table}[h]
\centering
\begin{tabular}{lcccc}
\hline
Metric & Trinity & Random (median) & $p$\-value \\
\hline
Mean rel.\ error & 0.0835\% & 0.594\% & $<10^{-4}$ \\
Hits $<0.1\%$ & 20/25 & 5/25 & $<10^{-4}$ \\
Hits $<1.0\%$ & 25/25 & 21/25 & 0.0061 \\
\hline
\end{tabular}
\caption{Monte\-Carlo comparison (50\,000 trials, 2\,000 formulas/trial).}
\label{tab:mc}
\end{table}

Table~\ref{tab:mc} shows that the \emph{collective} precision of the
catalog is highly significant ($p<10^{-4}$).  However, only 10 of the
25 formulas are individually significant at $p<0.05$; the remaining 15
are consistent with random search when tested in isolation.  The
strength of the catalog is therefore \emph{joint}, not individual.

A second test using real PDG 2024 uncertainties gives a median
$\sigma$\-distance of $0.12\sigma$ (half the formulas agree within
$1/8$ of the experimental uncertainty).  Three formulas fail
$\sigma$\-distance catastrophically: $1/\alpha$ ($1.6\times10^6\sigma$),
$m_p/m_e$ ($3.1\times10^5\sigma$) and $m_\mu/m_e$ ($6\times10^3\sigma$).
This is not because their relative errors are large (all $<0.03\%$),
but because these are the most precisely measured constants in physics;
any claim to ``derive'' them must match experimental precision.

\subsection{What survives?}

Despite the Boundary theorems, several structural observations survive:
\begin{itemize}
  \item The eta-invariant $-2$ is exact and physically meaningful (APS).
  \item KO-dim $6 \pmod 8$ for H$_4$ is a genuine geometric fact.
  \item The 25 verified formulas suggest \emph{accidental} (not derived)
    numerical coincidences, valuable as a ``periodic table'' of SM
    parameters.  The Monte\-Carlo test (Sec.~\ref{sec:stats}) confirms
    that this ``periodic table'' is statistically non\-trivial.
  \item JUNO 2025 measures $\Delta m^2_{21} = (7.50 \pm 0.12) \times 10^{-5}\,\text{eV}^2$
    \cite{JUNO:2025}, consistent with the Trinity value at $0.25\sigma$.
  \item DESI DR2 favors normal hierarchy with $\sum m_\nu < 0.064\,\text{eV}$
    \cite{DESI:2025}, consistent with Trinity's normal-hierarchy prediction.
\end{itemize}

\subsection{Experimental falsification registry}

The Trinity program maintains a pre-registered falsification tracker:
\begin{itemize}
  \item $m_H = 132.88\,\text{GeV}$ --- \textbf{REFUTED} by LHC (2012--2024)
  \item $\delta_{CP} = 65.66^\circ$ --- \textbf{>3$\sigma$ TENSION} with
    NuFit~6.0 ($\sim$$220^\circ \pm 30^\circ$). DUNE (2028--2032) will
    render a binary verdict.
  \item $\sin^2\theta_{13} = 0.0220$ --- \textbf{OPEN}, testable at JUNO
  \item $\Delta m^2_{21} = 7.53 \times 10^{-5}\,\text{eV}^2$ ---
    \textbf{CONSISTENT} with JUNO 2025 at $0.25\sigma$
  \item Normal mass hierarchy --- \textbf{CONSISTENT} with DESI DR2
\end{itemize}

\subsection{Relation to other programs}

Lisi's E$_8$ unification \cite{Lisi:2007} uses a different embedding
(fermions as E$_8$ roots). Our test shows that E$_8$ does not
automatically rescue the sigma-field either (Wave 13.3, open status).

The G\"{u}rsey--Ramond program \cite{Gursey:1975,Ramond:2001} proposed
D$_4$ triality as a source of three generations. Our F$_4$ scan
(\S\ref{sec:f4yukawa}) demonstrates that triality alone, without
symmetry-breaking parameters, cannot explain the observed mass hierarchy.
This strengthens the case that generation number requires a dynamical
mechanism beyond static Coxeter geometry.

String theory compactifications naturally obtain three generations from
intersection numbers, not from Coxeter-group triality. Our boundary finding
strengthens the case for string-derived, rather than discrete-geometric,
origin of generation number.

\subsection{Formal verification status}

Wave~16 achieved \textbf{zero admitted proofs} across the canonical proof
tree. The remaining unproven obligations are $77$ explicit Axioms, all tagged
as physical assumptions or mathematical gaps. One inline \texttt{admit}
remains in \texttt{E6vsH4.v}. The Lean~4 port now covers $9$ modules;
one honest \texttt{sorry} remains in \texttt{Spectrum600Cell.lean}
pending Mathlib integration.

\subsection{Open questions}

\begin{enumerate}
  \item Can a rigorously derived B-matrix for the E$_8$ plumbing improve
    the discrete eta convergence?
  \item Can the 1-loop Coleman--Weinberg potential on the 600-cell be
    modified to reach $125\,\text{GeV}$ without the $\sigma$-field?
  \item Are there experimental signatures (e.g.~axion-like $f_0$) that
    would distinguish the Trinity framework from vanilla SM?
\end{enumerate}

% ============================================================
\section{Conclusion}\label{sec:conclusion}

The Trinity S$^3$AI program set out to derive the Standard Model from the
H$_4$ Coxeter group and the 600-cell. After 17 waves of formal proof,
numerical validation, and honest refutation, we conclude:

\begin{quote}
  \textbf{The H$_4$/600-cell discrete model does not derive the Standard
  Model.} The sigma-field is absent, three generations are absent, the
  Higgs mass is off by $> 50\sigma$, and the F$_4$ extension fails the
  Yukawa hierarchy test.
\end{quote}

However, the program produced \emph{constructive} boundary theorems:
exact theorems, falsifiable predictions, and a rigorous framework for
testing future geometric unification proposals. We encourage the community
to use our open-source Coq/Lean/Python infrastructure to test alternative
models.

% ============================================================
\section*{Acknowledgements}

We thank the Coq, Lean, and NCG communities for their libraries and
documentation. This work was supported entirely by curiosity and
open-source tooling.

% ============================================================
\appendix

\section{Coq Proof Statistics}\label{app:coq}

As of Wave 17:
\begin{verbatim}
  Qed:        1375
  Admitted:      0
  Axioms:       77 (all tagged)
  Refutations:   7 (Boundary Theorems 1--5, 7)
  Files:        79 .v files
  CI:           passing (Coq 8.20.1)
\end{verbatim}

\section{Lean 4 Port Status}\label{app:lean}

\begin{itemize}
  \item \texttt{CorePhi.lean}: fully proved
  \item \texttt{KODimension.lean}: fully proved
  \item \texttt{QuaternionicLinearity.lean}: fully proved (11 Float axioms)
  \item \texttt{Spectrum600Cell.lean}: 1 honest \texttt{sorry} (needs Mathlib)
  \item \texttt{EtaInvariant.lean}: fully proved
  \item \texttt{DiracOperator.lean}: definitions solid
\end{itemize}

\section{Python Validation Results}\label{app:python}

\begin{itemize}
  \item \texttt{validate\_v4.py}: 25/25 PASS, mean error $0.10\%$
  \item \texttt{d4\_analysis.py}: D$_4$ spectrum computed, KO-dim $= 5$
  \item \texttt{spectral\_action.py}: Higgs $132.88\,\text{GeV}$
  \item \texttt{build\_e6\_e7\_dp.py}: plumbing operators constructed
  \item \texttt{f4\_spectrum.py}: F$_4$ spectrum, KO-dim $= 6$
  \item \texttt{yukawa\_scan.py}: F$_4$ Yukawa scan --- symmetric models FAIL
  \item \texttt{b\_matrix\_derivation.py}: E$_8$ plumbing $D_P$ (512$\times$512)
\end{itemize}

\section{Falsification Registry}\label{app:falsify}

Key predictions and their experimental status:
\begin{itemize}
  \item $m_H = 132.88\,\text{GeV}$ --- REFUTED by LHC (2012--2024)
  \item $\delta_{CP} = 65.66^\circ$ --- $>$$3\sigma$ TENSION with NuFit 6.0;
    DUNE verdict 2028--2032
  \item $\sin^2\theta_{13} = 0.0220$ --- OPEN, testable at JUNO/DUNE
  \item $\Delta m^2_{21} = 7.53 \times 10^{-5}\,\text{eV}^2$ --- CONSISTENT
    with JUNO 2025
  \item Normal hierarchy --- CONSISTENT with DESI DR2
\end{itemize}
Full registry in \texttt{PREDICTIONS\_REGISTRY.md}.

% ============================================================
\begin{thebibliography}{99}

\bibitem{Connes:1996}
A. Connes, ``Gravity coupled with matter and the foundation of
non-commutative geometry,'' \textit{Commun. Math. Phys.} \textbf{182} (1996)
155--176.

\bibitem{ConnesMarcolli:2008}
A. Connes and M. Marcolli, \textit{Noncommutative Geometry, Quantum Fields
and Motives}, AMS, 2008.

\bibitem{Lisi:2007}
G. Lisi, ``An exceptionally simple theory of everything,''
\textit{arXiv:0711.0770 [hep-th]}.

\bibitem{ConwaySmith:2003}
J. H. Conway and D. A. Smith, \textit{On Quaternions and Octonions},
AK Peters, 2003.

\bibitem{ElserSloane:1987}
V. Elser and N. J. A. Sloane, ``A highly symmetric 4-D tiling,''
\textit{J. Phys. A} \textbf{20} (1987) 6161--6168.

\bibitem{APS:1975}
M. Atiyah, V. Patodi, and I. Singer, ``Spectral asymmetry and
Riemannian geometry,'' \textit{Math. Proc. Camb. Phil. Soc.} \textbf{77}
(1975) 43--69.

\bibitem{PDG:2024}
Particle Data Group, \textit{Review of Particle Physics}, Phys.\ Rev.\ D
\textbf{110}, 030001 (2024).

\bibitem{ChamseddineConnes:1997}
A. H. Chamseddine and A. Connes, ``The spectral action principle,''
\textit{Commun. Math. Phys.} \textbf{186} (1997) 731--750.

\bibitem{Gilkey:1984}
P. B. Gilkey, \textit{Invariance Theory, the Heat Equation, and the
Atiyah--Singer Index Theorem}, Publish or Perish, 1984.

\bibitem{Koca:2006}
M. Koca et al., ``Coxeter groups and the H4 polytope,''
\textit{J. Phys. A} \textbf{39} (2006) 7755--7780.

\bibitem{JUNO:2025}
JUNO Collaboration, ``First physics results from the Jiangmen Underground
Neutrino Observatory,'' \textit{arXiv:2511.14593 [hep-ex]} (2025).

\bibitem{DESI:2025}
DESI Collaboration, ``DESI 2024 VI: Cosmological constraints from the
measurements of baryon acoustic oscillations,''
\textit{Phys.\ Rev.\ D} \textbf{112}, 083515 (2025).

\bibitem{Ramond:2001}
P. Ramond, ``Exceptional groups and physics,''
\textit{arXiv:hep-th/0112261}.

\bibitem{Gursey:1975}
F. G\"{u}rsey, P. Ramond, and P. Sikivie,
``A universal gauge theory model based on E$_6$,''
\textit{Phys.\ Lett.\ B} \textbf{60} (1975) 177.

\end{thebibliography}

\end{document}
